\chapter{Multiplicity and Sequentially Rejective Testing}
\label{sec:p15-multiple-testing}

A family of individually valid tests can still make a false rejection
with substantial probability. The relevant error event is that at least
one hypothesis that is true is rejected. Holm's sequentially rejective
Bonferroni procedure controls this event for every configuration of true
hypotheses, while allowing thresholds to increase after earlier
rejections \cite[pp.~65--70]{Holm1979Original}.
Its sequential steps inspect one fixed set of test results.
They do not authorize repeated sampling until significance appears.

\section{The family, its sampling law and obtained levels}
\label{sec:p15-mt-family}

\begin{hbdefinition}{def:p15-mt-family}{True hypotheses and the familywise error event}
Use the statistical experiment
\((\mathsf X,\mathscr X,\Theta,(P_\theta)_{\theta\in\Theta})\)
of \cref{def:p15-statistical-experiment}. Fix a finite family
\(H_1,\ldots,H_n\), \(n\geq1\), before observing \(X\).
The hypothesis \(H_i\) asserts \(\theta\in\omega_i\), where
\(\varnothing\ne\omega_i\subseteq\Theta\).
For a given parameter define
\[
 I_0(\theta)=\{i:\theta\in\omega_i\},\qquad m=\abs{I_0(\theta)}.
\]
These are the actual true-null indices and their number; they are not
observed inputs to the procedure.

A multiple test is a measurable rejection set
\(\mathcal R(X)\subseteq\{1,\ldots,n\}\): each event
\(\{i\in\mathcal R(X)\}\) is measurable. Its critical region for \(H_i\)
is \(C_i=\{x:i\in\mathcal R(x)\}\), and its familywise error probability is
\[
 \operatorname{FWER}_\theta(\mathcal R)
   =P_\theta\{\mathcal R(X)\cap I_0(\theta)\ne\varnothing\}.
\]
Strong control at level \(\alpha\), \(0<\alpha<1\), means this probability
is at most \(\alpha\) for every \(\theta\).
When \(m=0\) the event is empty. Control only when all hypotheses are true
is a weaker property. Nonrejection does not prove a hypothesis.
\index{familywise error}\index{multiple testing}
\end{hbdefinition}

\begin{hbdefinition}{def:p15-mt-obtained-levels}{Marginal validity and Holm's critical-level construction}
For each \(i\), let \(R_i:\mathsf X\to[0,1]\) be measurable and suppose
\[
 P_\theta\{R_i(X)\leq u\}\leq u
       \qquad(\theta\in\omega_i,\ 0\leq u\leq1).                 \tag{V}
\]
Thus \(R_i\) is a valid p-value in the sense of
\cref{def:p15-valid-p-value}; Holm calls it an obtained level.
The inequalities hold under the full sampling parameter, including its
nuisance parameters and the truth or falsity of the other hypotheses.
There is no assumption about the joint dependence of the \(R_i\).

For a real-valued statistic \(Y_i\), with large values favoring rejection,
Holm's construction is
\[
 \alpha_i(y)=\sup_{\vartheta\in\omega_i}P_\vartheta(Y_i\geq y),
 \qquad R_i=\alpha_i(Y_i).
\]
Here \(\alpha_i(y)\) is a tail function, distinct from the chosen constant
family level \(\alpha\). Inclusive tails retain their validity at atoms.
The supremum tail is nonincreasing and hence Borel measurable.
An approximate tail calculation needs its own calibration error bound;
calling a value a p-value does not imply (V).
\end{hbdefinition}

\begin{hblemma}{thm:p15-mt-inclusive-levels}{proved=annex:p15.mt.inclusive-levels}{Inclusive composite-null tails are valid}
For any real statistic \(Y_i\), the displayed supremum-tail construction
satisfies (V), including discrete distributions, atoms and composite
hypotheses. No continuity, identical marginal distribution across tests,
or independence across tests is required.
\end{hblemma}

\section{The step-down rule and its protection}
\label{sec:p15-mt-step-down}

\begin{hbmethod}{def:p15-mt-holm-rule}{The exact sequentially rejective Bonferroni procedure}
At the observed data order the obtained levels as
\[
 R^{(1)}\leq R^{(2)}\leq\cdots\leq R^{(n)},\qquad
 R^{(j)}=R_{i_j},\quad H^{(j)}=H_{i_j}.
\]
Break ties by a fixed ordering of the original indices.
Starting with \(j=1\), reject \(H^{(j)}\) if
\[
                    R^{(j)}\leq\frac{\alpha}{n-j+1}.
\]
If this comparison fails, stop and do not reject any remaining hypothesis.
If it passes and \(j<n\), increase \(j\) by one; if \(j=n\), stop.
Denote the resulting rejection set by \(\mathcal R_{\rm H}\).
Equivalently, \(H^{(r)}\) is rejected exactly when every comparison for
\(1\leq j\leq r\) passes. Equality passes, as in Holm's Scheme~1.
The procedure makes no decisions beyond this specified family.
\end{hbmethod}

\begin{hblemma}{thm:p15-mt-ordering}{proved=annex:p15.mt.ordering}{Measurability, ties and termination}
The rule terminates after at most \(n\) comparisons and is a measurable
multiple test. Changing the order within a tied block does not change
which hypotheses are rejected: a tied block is rejected in its entirety
or not at all. For \(n=1\), the rule is the single test \(R_1\leq\alpha\).
The empty-family convention, if needed, returns the empty rejection set.
\end{hblemma}

\begin{hbtheorem}{thm:p15-mt-holm-strong}{proved=annex:p15.mt.holm-strong}{Holm's protection for every configuration}
Under (V),
\[
                  \operatorname{FWER}_\theta(\mathcal R_{\rm H})
                         \leq\alpha\qquad(\theta\in\Theta).
\]
For \(m\geq1\), a false rejection implies
\[
             \min_{i\in I_0(\theta)}R_i\leq\alpha/m.
\]
Consequently the result follows from the marginal bounds and the finite
union inequality, without an independence assumption.

Equivalently, for every nonempty \(I\subseteq\{1,\ldots,n\}\) and every
\(\theta\in\bigcap_{i\in I}\omega_i\), the probability of rejecting any
hypothesis indexed by \(I\) is at most \(\alpha\). This is Holm's
\emph{free combinations} protection. Restrictions on which truth
configurations can occur do not invalidate it.
\end{hbtheorem}

\section{Comparisons and adjusted levels}
\label{sec:p15-mt-comparisons}

\begin{hbdefinition}{def:p15-mt-comparison-rules}{Fixed Bonferroni and unadjusted rejection sets}
Using the same family and same obtained levels, put
\[
 \mathcal R_{\rm B}(x)=\{i:R_i(x)\leq\alpha/n\},\qquad
 \mathcal R_{\rm U}(x)=\{i:R_i(x)\leq\alpha\}.
\]
The first is the fixed Bonferroni rule; the second applies the nominal
single-test level to every member of the family.
\end{hbdefinition}

\begin{hbtheorem}{thm:p15-mt-domination}{proved=annex:p15.mt.domination}{Pointwise dominance and its precise power meaning}
For every observation,
\[
               \mathcal R_{\rm B}\subseteq
               \mathcal R_{\rm H}\subseteq\mathcal R_{\rm U}.
\]
The Bonferroni rule also has strong FWER control under (V).
For any fixed set \(J\) of false hypotheses at a specified parameter,
Holm's probability of rejecting all of \(J\), its probability of rejecting
at least one of \(J\), and its expected number of rejections in \(J\)
are each at least the corresponding Bonferroni quantity.
These inequalities need not be strict.
Holm's probability of making no rejection is exactly Bonferroni's.
\end{hbtheorem}

\begin{hbdefinition}{def:p15-mt-adjusted}{Adjusted levels for the whole step-down decision}
In the observed ordering define
\[
 \widetilde R^{(r)}
  =\min\left\{1,\max_{1\leq j\leq r}(n-j+1)R^{(j)}\right\}.
\]
Assign this adjusted level to the original hypothesis \(H_{i_r}\).
The maximum retains every earlier stopping condition; it is not enough
to multiply each p-value by its own remaining-family count alone.
\end{hbdefinition}

\begin{hblemma}{thm:p15-mt-adjusted}{proved=annex:p15.mt.adjusted}{Exact adjusted-level equivalence}
For \(0<\alpha<1\), Holm rejects \(H_{i_r}\) if and only if
\(\widetilde R^{(r)}\leq\alpha\). The adjusted levels are nondecreasing
in rank and equal throughout a tied block. Thus both the decisions and
these levels are independent of the choice of tie ordering.
\end{hblemma}

\begin{hbexample}{def:p15-mt-worked-family}{A stopping condition that later comparisons cannot repair}
For five hypotheses at \(\alpha=0.05\), suppose the sorted obtained levels
are \(0.004,0.012,0.018,0.019,0.040\). Then
\[
\begin{array}{c|ccccc}
 j&1&2&3&4&5\\ \hline
 R^{(j)}&0.004&0.012&0.018&0.019&0.040\\
 \alpha/(6-j)&0.010&0.0125&0.016\overline6&0.025&0.050\\
 (6-j)R^{(j)}&0.020&0.048&0.054&0.038&0.040\\
 \widetilde R^{(j)}&0.020&0.048&0.054&0.054&0.054
\end{array}
\]
Holm rejects the first two hypotheses and stops at the third. The fourth
and fifth hypothetical comparisons would pass in isolation, but they
are never reached. Bonferroni rejects only the first; unadjusted testing
rejects all five and supplies no strong family guarantee.
\end{hbexample}

\section{Intersection hypotheses, direction and confidence sets}
\label{sec:p15-mt-logical-structure}

\begin{hbdefinition}{def:p15-mt-intersection-policy}{The source's secondary hypotheses}
For a nonempty index set \(I\), the secondary hypothesis \(H_I\) asserts
\(\theta\in\bigcap_{i\in I}\omega_i\).
Given the basic rejection set \(\mathcal R\), declare \(H_I\) rejected
when \(I\cap\mathcal R\ne\varnothing\).
For \(I\subseteq J\), \(H_J\) implies \(H_I\).
Coherence for these containment implications means that rejection of
\(H_I\) also rejects \(H_J\). Consonance means that rejection of \(H_I\)
has a rejected basic member \(H_i\), \(i\in I\).
Additional semantic implications among the basic hypotheses require
their own consistency analysis; containment is not an inventory of all
possible logical relations.
\end{hbdefinition}

\begin{hblemma}{thm:p15-mt-intersection-policy}{proved=annex:p15.mt.intersection-policy}{Coherence and consonance of the stated secondary policy}
The policy in \cref{def:p15-mt-intersection-policy} is coherent for
containment implications and consonant. Under strong basic FWER control,
the probability of falsely rejecting any true secondary hypothesis is
also at most \(\alpha\).
\end{hblemma}

\begin{hbtheorem}{thm:p15-mt-closed-bonferroni}{proved=annex:p15.mt.closed-bonferroni}{A complete closed-testing realization}
For each nonempty \(I\), use the local intersection test that rejects
\(H_I\) when
\[
                         \min_{i\in I}R_i\leq\alpha/\abs I.
\]
This local test has level at most \(\alpha\) on the intersection null.
Reject a basic \(H_i\) precisely when every intersection test whose index
set contains \(i\) rejects. The resulting basic rejection set is exactly
\(\mathcal R_{\rm H}\). The equivalence proves the relation to closed
testing without replacing the explicit step-down construction.
\end{hbtheorem}

\begin{hblemma}{thm:p15-mt-directional}{proved=annex:p15.mt.directional}{Opposite tails avoid contradictory directional conclusions}
Let two family members be \(H_a:\gamma\leq\gamma_0\) and
\(H_b:\gamma\geq\gamma_0\), with a common boundary law
at \(\gamma=\gamma_0\) present in the sampling model.
Use statistics \(Y\) and \(-Y\) and the
inclusive supremum tails of \cref{def:p15-mt-obtained-levels}.
Then \(R_a(y)+R_b(-y)\geq1\) at every observation.
For \(0<\alpha\leq1/2\), the ordinary Holm procedure cannot reject both.
Thus a family of such pairs avoids opposite directional rejections when
these are its only additional consistency constraints.
FWER protection alone does not assert consistency for arbitrary other
logical relations.
\end{hblemma}

\begin{hblemma}{thm:p15-mt-confidence-identical}{proved=annex:p15.mt.confidence-identical}{The source's confidence-set limitation}
At each candidate parameter \(v\), suppose jointly measurable obtained
levels \(R_i(x;v)\) are used to form the set of candidates for which
no detail hypothesis is rejected. The set obtained with ordinary Holm
is exactly
\[
 \mathcal C(x)=\{v:R_i(x;v)>\alpha/n\text{ for every }i\},
\]
which is the set obtained with fixed Bonferroni.
If all detail hypotheses are true at the actual candidate and their
p-values satisfy (V), this set has coverage at least \(1-\alpha\).
The equality concerns this all-details-nonrejection inversion. It is
not a statement about every possible simultaneous interval construction.
\end{hblemma}

\section{The original power comparison}
\label{sec:p15-mt-normal-power}

\begin{hbdefinition}{def:p15-mt-normal-power}{Ten independent normal tests}
Let the \(Y_i\) be independent \(N(\mu_i,1)\) variables. Test
\(H_i:\mu_i=0\) against \(\mu_i>0\) with
\(R_i=1-\Phi(Y_i)\), where \(\Phi\) is the standard normal CDF.
At \(\alpha=0.05\), take four means equal to zero, four equal to six,
and two equal to three, as in Holm's example
\cite[pp.~67--68]{Holm1979Original}.
The power quantity here is the probability that both hypotheses with
mean three are rejected. Independence belongs to this example, not to
the general Holm guarantee.
\end{hbdefinition}

\begin{hblemma}{thm:p15-mt-normal-power}{proved=annex:p15.mt.normal-power}{A finite formula for the two designated rejections}
Put \(b_0=0\), \(b_k=0.05/(11-k)\), \(1\leq k\leq10\).
For a mean \(\mu\), the obtained-level CDF is
\[
 F_\mu(u)=\Phi\bigl(\mu-\Phi^{-1}(1-u)\bigr)\quad(0<u<1),
 \qquad F_\mu(0)=0,\quad F_\mu(1)=1.
\]
Partition levels into the bins \((b_{k-1},b_k]\) and \((b_{10},1]\).
After all earlier comparisons have passed, if \(a\) levels have occurred
in earlier bins, the first level in occupied bin \(k\) passes Holm
exactly when \(a+1\geq k\).
Every other level in that bin then passes.
These finite bin counts determine whether both designated mean-three
tests reject; the annex gives the full probability recursion.

Using exact normal quantiles, Bonferroni's probability is
\[
 \Phi\bigl(3-\Phi^{-1}(0.995)\bigr)^2
                  \simeq0.4412670349,
\]
and evaluation of the finite Holm recursion gives approximately
\(0.5631999191\). These are numerical evaluations of the displayed
sampling model; the increased power does not remove its modeling assumptions.
\end{hblemma}

\section{The original weighted and independent extensions}
\label{sec:p15-mt-extensions}

\begin{hbmethod}{def:p15-mt-weighted}{Positive fixed weights and Scheme 2}
Choose \(c_1,\ldots,c_n>0\) independently of the observed data.
Form \(S_i=R_i/c_i\), sort
\(S^{(1)}\leq\cdots\leq S^{(n)}\), and carry the corresponding
\(i_j,H^{(j)}\) and \(c^{(j)}\) with them, with a fixed tie order.
These indices and superscripts now refer to the ratio ordering.
At step \(j\), put \(W_j=\sum_{\ell=j}^n c^{(\ell)}\).
Reject the current hypothesis if
\[
                  S^{(j)}\leq\alpha/W_j,
       \quad\text{equivalently}\quad
                  R_{i_j}\leq\alpha c^{(j)}/W_j.
\]
Stop at the first failure. Denote the rejection set by
\(\mathcal R_{\rm W}\). The remaining comparison budget is divided in
proportion to the remaining positive constants, as in Holm's Scheme~2
\cite[pp.~69--70]{Holm1979Original}.
The ordering uses the ratios, not the unweighted levels.
\end{hbmethod}

\begin{hbtheorem}{thm:p15-mt-weighted}{proved=annex:p15.mt.weighted}{Full strong protection for the weighted procedure}
Under (V), \(\mathcal R_{\rm W}\) has strong FWER control at \(\alpha\).
If \(C_0=\sum_{i\in I_0(\theta)}c_i>0\), its first true rejection implies
\(R_i\leq\alpha c_i/C_0\) for some true \(i\).
The guarantee follows by summing these true-hypothesis allocations.
Common rescaling of all weights does not change the rule; equal weights
recover ordinary Holm. Ratio ties yield whole-block decisions and do
not depend on the tie order. The weighted procedure contains the
fixed weighted Bonferroni set
\(\{i:R_i\leq\alpha c_i/\sum_\ell c_\ell\}\).
No general dominance over unweighted Holm is asserted.
Under the opposite-tail conditions of \cref{thm:p15-mt-directional},
it likewise cannot reject both members when \(\alpha\leq1/2\).
\end{hbtheorem}

\begin{hbexample}{def:p15-mt-weight-tradeoff}{What weighting can change}
Take three hypotheses with \((c_1,c_2,c_3)=(8,1,1)\) and \(\alpha=0.05\).
For obtained levels \((0.030,0.040,0.040)\), unweighted Holm makes no
rejection, whereas the weighted procedure rejects \(H_1\) and then stops.
For \((0.048,0.008,0.009)\), unweighted Holm rejects all three.
The weighted ordering puts \(H_1\) first, its first threshold is \(0.040\),
and it stops without rejecting anything. The chosen weights specify
priority; their benefits and losses depend on the alternative configuration.
\end{hbexample}

\begin{hbmethod}{def:p15-mt-independent}{The independence-based step-down extension}
In the ordinary obtained-level ordering, replace \(\alpha/(n-j+1)\) by
\[
                  d_{n-j+1}=1-(1-\alpha)^{1/(n-j+1)}
\]
and keep the same stopping rule.
This is the independent-tail extension described by Holm on p.~68,
often called step-down \v{S}id\'ak testing.
Its justification requires the true obtained levels to be mutually
independent under every parameter at which those hypotheses are true.
Independence solely at the complete null is insufficient for strong control.
\end{hbmethod}

\begin{hbtheorem}{thm:p15-mt-independent}{proved=annex:p15.mt.independent}{The exact independence premise and the larger thresholds}
Assume (V) and, at each \(\theta\), mutual independence of
\((R_i)_{i\in I_0(\theta)}\). The independent-tail procedure has strong
FWER control. It does not require independence between a false-hypothesis
level and a true-hypothesis level.
For \(k\geq1\), \(d_k\geq\alpha/k\), with strict inequality for \(k>1\).
Its rejection set therefore contains ordinary Holm's pointwise.
Under arbitrary dependence the protection can fail.
\end{hbtheorem}

\section{Applications with the actual marginal assumptions}
\label{sec:p15-mt-applications}

Holm's application discussion covers comparisons of treatments with one
control, unequal sample sizes, directional comparisons, nonparametric
many-to-one comparisons, contingency tables, time series and multivariate
distributions \cite[pp.~68--70]{Holm1979Original}.
The family rule remains the same; the required marginal calibration
must be established for the actual sampling model.

\begin{hbdefinition}{def:p15-mt-control-model}{Normal comparisons with a common control}
There are \(n\) treatments and a control indexed by \(g=0\).
Group \(g\) supplies \(N_g\geq1\) independent normal observations
with mean \(\mu_g\) and common variance \(\sigma^2>0\); groups are independent.
Put \(N_{\rm tot}=\sum_{g=0}^nN_g\) and require
\(\nu=N_{\rm tot}-(n+1)>0\). Define
\[
 s_{\rm p}^2=\frac1\nu\sum_{g=0}^n\sum_{\ell=1}^{N_g}
                     (X_{g\ell}-\bar X_g)^2,\qquad
 T_i=\frac{\bar X_i-\bar X_0}
                   {s_{\rm p}\sqrt{1/N_i+1/N_0}}\quad(s_{\rm p}>0).
\]
Test \(H_i:\mu_i\leq\mu_0\) with \(R_i=1-F_\nu(T_i)\), using the
Student CDF of \cref{thm:p15-student-law}. On the measurable zero-scale
set define every \(R_i=1\); this set has probability zero under the model.
The procedure is defined on every sample outcome.
\end{hbdefinition}

\begin{hblemma}{thm:p15-mt-control-model}{proved=annex:p15.mt.control-model}{Valid unequal-size comparisons despite shared control and variance}
Each displayed \(R_i\) satisfies (V) under its entire one-sided null.
Consequently Holm, including its weighted form, controls the family error
despite the shared control and pooled variance.
For \(n=9\) and \(N_g=4\), \(\nu=30\); the Holm critical Student values
at family level \(0.05\) are \(t_{30,1-0.05/k}\), where \(k\) is the
number of remaining comparisons.
\end{hblemma}

For that equal-size example the source prints the following rounded
critical values. Its refined Dunnett values are reported for its
particular comparison model; they are not universal Holm thresholds:
\[
\begin{array}{c|rrrrrrrrr}
 k&1&2&3&4&5&6&7&8&9\\ \hline
 \text{refined Dunnett}&1.70&1.99&2.15&2.25&2.33&2.40&2.45&2.50&2.54\\
 \text{Holm}&1.70&2.04&2.23&2.36&2.46&2.54&2.60&2.66&2.71
\end{array}
\]
The first Holm comparison is the largest statistic against the \(k=9\)
critical value, then \(k=8\), and so on. Classical Dunnett compares
each statistic with \(2.54\) in the source's example.
These model-specific comparisons do not imply that Holm dominates
procedures that use the joint model more efficiently.
With unequal \(N_g\), the marginal Student calculation above still
applies under its common-variance assumptions. Heterogeneous variances
require differently justified marginal tests.

\begin{hblemma}{thm:p15-mt-conditional-calibration}{proved=annex:p15.mt.conditional-calibration}{Conditional tests supply valid marginal levels}
If a test level \(R_i\) obeys
\[
          P_\theta(R_i\leq u\mid Z_i)\leq u\quad\text{almost surely}
\]
for every \(\theta\in\omega_i\), then it obeys (V).
For a finite conditional permutation distribution that is uniform under
the null, the inclusive fraction of permutations at least as extreme
as the observed statistic has this property, including ties.
Different tests may condition on different information; no independence
among their obtained levels is needed.
\end{hblemma}

For nonnormal treatment comparisons, the source proposes ordinary
two-sample Wilcoxon tests as the marginal inputs. Under the null,
exchangeability of the pooled observations gives the finite permutation
law, so the preceding conditional calibration applies.
An arbitrary difference-of-means interpretation without that null
exchangeability is not supplied by the rank test.
For contingency tables, use marginal tests whose exact conditional or
unconditional null laws match the table model. For time series and
multivariate data, dependence among the tests is permitted, but dependence
within each statistic's sampling model still affects its validity.
The source's comments on table availability concern its 1979 computing
setting; they impose no restriction on which joint procedures can now
be computed.

\section{Failure cases and calibration boundaries}
\label{sec:p15-mt-failures}

\begin{hbtheorem}{thm:p15-mt-failures}{proved=annex:p15.mt.failures}{Four explicit failures of neighboring procedures}
The following constructions distinguish the theorem's hypotheses.
\begin{enumerate}
\item For twenty independent true-null levels uniform on \((0,1)\),
unadjusted \(0.05\) testing has FWER
\(1-0.95^{20}\simeq0.6415141\).
Holm and Bonferroni have the same any-rejection event at this complete
null and FWER \(1-(1-0.05/20)^{20}\simeq0.0488301\).
Choosing only the smallest observed level and calling it a family of one
recreates the unadjusted error.
\item Rejecting the whole family if its smallest level is at most
\(\alpha/2\) has complete-null control for two independent uniform levels.
With one false hypothesis having level zero and the other a true
uniform level, it has FWER one. Complete-null control has not proved
the required every-configuration guarantee.
\item Two true-null levels \(U\) and \(1-U\) are marginally uniform.
The independent-tail procedure at \(\alpha=0.05\) has FWER
\(2(1-\sqrt{0.95})\simeq0.0506411>0.05\).
Ordinary Holm has FWER \(0.05\) in the same dependent model.
\item With two independent true-null uniform levels, choosing weights
99 and 1 after observing which level is smaller gives it a first-step
threshold \(99\alpha/100\). At \(\alpha=0.05\), its FWER is
\(1-(1-0.0495)^2=0.09654975\).
Marginal validity does not justify these data-selected weights.
\end{enumerate}
\end{hbtheorem}

\begin{hblemma}{thm:p15-mt-calibration-error}{proved=annex:p15.mt.calibration-error}{An explicit bound when validity is approximate}
Suppose that instead of (V), at a fixed parameter each true \(i\) has
\[
 P_\theta(R_i\leq u)\leq u+\varepsilon_i\quad(0\leq u\leq1),
 \qquad \varepsilon_i\geq0.
\]
Then ordinary and fixed-weight Holm both satisfy
\[
 \operatorname{FWER}_\theta
       \leq\min\left\{1,\alpha+\sum_{i\in I_0(\theta)}\varepsilon_i\right\}.
\]
Exact strong control therefore needs the actual marginal validity premise
or a calibration guarantee that restores it. It does not follow merely
from an asymptotic reference distribution.
\end{hblemma}

The remedy for a selected-family error is to define the intended family
and preserve its valid marginal levels through selection, or to supply
a separately justified selective procedure. The remedy for shared-data
dependence is the union-bound Holm rule when only marginal validity is
known. The remedy for a misspecified marginal sampling law is to repair
that calibration; changing the multiplicity threshold does not prove
the missing individual validity.
